\documentclass[a4paper]{article}

\usepackage[margin=1in]{geometry}
\usepackage{booktabs}
\usepackage{graphicx}
\usepackage{parskip}
\usepackage{hyperref}

\usepackage{libertinust1math}
\usepackage{tgpagella}

\title{Causality Course}
\author{Dan Andrei Iliescu}
\date{May 2022 --- May 2023}

\begin{document}

\maketitle

\begin{figure}[h]
    \includegraphics[width=\textwidth]{files/2023_03_24.pdf}
    \caption{\textbf{The Potential Outcomes Framework.}}
\end{figure}

\begin{figure}[h]
    \includegraphics[width=\textwidth]{files/2023_04_28.pdf}
    \caption{\textbf{The Backdoor Criterion.}}
\end{figure}

\begin{figure}[h]
    \includegraphics[width=\textwidth]{files/2023_05_26.pdf}
    \caption{\textbf{Covariate Adjustment and Frontdoor Adjustment.}}
\end{figure}

\section{Frontdoor Adjustment --- 2023--06--01}

\paragraph{Theorem 6.1} If $(T, M, Y)$: 

\begin{enumerate}
    \item Satisfy the frontdoor criterion.
    \item We have positivity. Positivity means that $0 < P(t | x) < 1, ~ \forall x, t$.
\end{enumerate}

then 

\begin{equation}
    P(y | \mathrm{do} (t)) = \sum_{m} P(m | t) \sum_{t'} P(y | m, t') P(t')    
\end{equation}

So what is the frontdoor criterion?

\paragraph{Definition 6.1} A set of variables $M$ satisfy the frontdoor criterion relative to $T, Y$ if:

\begin{enumerate}
    \item All causal paths $T \to Y$ pass through $M$.
    \item There is no unblocked backdoor path $T \to M$.
    \item All backdoor paths $M \to Y$ are blocked when conditioning on $T$.
\end{enumerate}

\paragraph{Reminder.} A backdoor path $A \to B$ is a non-direct causal association from $A$ to $B$. In my interpretation, a backdoor path is any chain of variables that connect $A$ to $B$ and that make $P(B | A) \neq P(B | \mathrm{do}(A))'$. 

\paragraph{Do-Calculus} The rules of do-calculus:

\begin{enumerate}
    \item Rule 1: $P(y | do(t), z, w) = P(y | do(t), w)$, if $Y \perp_{\overline{T}} Z | T, W$.
    \item Rule 2: $P(y | do(t), do(z), w) = P(y | do(t), z, w)$, if $Y \perp_{\overline{T}, \underline{Z}} Z | T, W$.
    \item Rule 3: $P(y | do(t), do(z), w) = P(y | do(t), w)$, if $Y \perp_{\overline{T}, \overline{Z(W)}} Z | T, W$.
\end{enumerate}

\paragraph{Proof 6.1} First, let's establish what $P(m | \mathrm{do}(t))$ and $P(y | \mathrm{do}(m))$ are.

\begin{enumerate}
    \item What is $P(m | \mathrm{do}(t))$? Are there any backdoors $T \to M$? No, therefore $P(m | \mathrm{do}(t)) = P(m | t)$.
    \item What is $P(y | \mathrm{do}(m))$? Are there any backdoors $M \to Y$? Yes, $M \leftarrow T \leftarrow W \rightarrow Y$. This path can be blocked by $T$, which is observed. Therefore $P(y | \mathrm{do}(m)) = \sum_{t} P(y, t | m) = \sum_{t} P(y | t, m) P(t)$.
\end{enumerate}

At this point, we know that

\begin{equation}
    \sum_{m} P(m | t) \sum_{t'} P(y | m, t') ~ P(t') = \sum_{m} P(m | \mathrm{do} (t)) ~ P(y | \mathrm{do} (m))
\end{equation}

Also, by the law of total probability, we know that

\begin{equation}
    P(y | \mathrm{do} (t)) = \sum_{m} P(y, m | \mathrm{do} (t)) = \sum_m P(y | m, \mathrm{do} (t)) ~ P(m | \mathrm{do} (t))
\end{equation}

So all we have to do is prove that

\begin{equation}
    \sum_m P(y | m, \mathrm{do} (t)) ~ P(m | \mathrm{do} (t)) = \sum_{m} P(y | \mathrm{do} (m)) ~ P(m | \mathrm{do} (t))
    \label{eq:frontdoor}
\end{equation}

We will now prove that $P(y | m, \mathrm{do} (t)) = P(y | \mathrm{do} (m))$. If we do that, then we're done. We apply the first two rules of do-calculuts to the termo on the left.

\begin{align}
    P(y | m, \mathrm{do} (t)) &= P(y | m, \mathrm{do} (t), \mathrm{do} (m)) \\
    &= P(y | \mathrm{do} (m))
\end{align}

And we're done! Therefore Equation~\ref{eq:frontdoor} is true.

\end{document}